\subsubsection{先判断是否需要 Conda}
一个 Python 项目要求这个版本的包，另一个项目又要求那个版本，如果都装进同一套环境里，很容易互相打架。这时候我们可以用环境管理器（Environment Manager），给每个项目准备一套相对独立的解释器和依赖，用哪个项目就切到哪套环境，不必每次都拆东墙补西墙。

Conda 就是同时管理解释器、包和环境的工具。如果你接触的 Python 项目不多，只需要系统 Python 和几个包，也可以跳过 Conda，用标准库自带的 venv，更轻量一些。
\begin{note}
环境隔离不是安全沙箱，也不保证所有依赖冲突都消失。安装器与下载渠道按当前官方说明核对：\url{https://docs.conda.io/projects/conda/en/stable/user-guide/install/}。已装好就不重复安装；\texttt{conda --version} 应能显示版本，\texttt{conda info --base} 能找到安装位置。
\end{note}

\subsubsection{先读配置，再决定是否更改}
安装好以后，我们先看看它现在怎样工作，不急着把网上的配置整段复制进去。Conda 可以从用户的 \texttt{\textasciitilde/.condarc} 等位置读取配置，\path|conda config --show-sources| 告诉你配置来自哪里，\texttt{--show 键名} 显示实际值，\texttt{--set 键名 值} 则用来设置单值。

下载慢时可以按当前渠道文档选择镜像，不过混用渠道可能影响依赖解算和许可适用条件，要按项目需要核对，不能把“加得越多就越容易装好”当成经验。
\Needspace{8\baselineskip}
\begin{example}[检查自动激活与下载渠道]


\Needspace{8\baselineskip}
\begin{tbash}
conda config --show-sources
conda config --show channels
conda config --show auto_activate_base
conda env list
\end{tbash}


先记下实际的渠道和环境列表。如果每次打开终端都会自动进入 base，而你并不总要用它，可以关闭基础环境的自动激活。已有 .condarc 时先备份自己的文件，然后执行：


\Needspace{5\baselineskip}
\begin{tbash}
conda config --set auto_activate_base false
\end{tbash}


接着从平时的入口新开一个终端，看看是不是不再自动进入 base。后面学习 ROS 2 时，Conda 的 PATH 或库路径可能干扰系统工具链，所以养成按需激活的习惯会更省心。如果你的版本提示键名或别名有变化，就用 \texttt{conda config --describe} 查当前支持项，不用猜一个名字反复改。
\end{example}
\begin{exercise}\label{ex:v1-conda-config}
只读记录配置来源、下载渠道与基础环境自动激活状态。说明哪些值来自用户配置、哪些来自默认值，以及为什么后续 ROS 2 编译前要确认没有继承 Conda 的 PATH/库路径。没有 Conda 时先使用 venv 的独立环境路线，Conda 命令检查留待安装后验证。
\end{exercise}

\subsubsection{创建、使用与删除练习环境}
配置看过了，我们来给一个小项目准备环境。创建时使用 \texttt{conda create --name 名称 python=版本 [包名]}，其中名称和版本按你的项目需要填写；创建好后，用 \texttt{conda activate 名称} 进入，\texttt{conda install 包名} 安装依赖，不用了再用 \texttt{conda deactivate} 退出。激活改变的是当前 Shell 使用的命令与依赖，旁边已经打开的其他终端不会自动跟着切换。
\Needspace{8\baselineskip}
\begin{example}[只给一个练习环境装 NumPy]
先在 \texttt{conda env list} 确认没有同名资料，再执行：


\Needspace{11\baselineskip}
\begin{tbash}
conda create --name rm-training-py310 python=3.10
conda activate rm-training-py310
conda install numpy
python --version
python -c "import sys, numpy; print(sys.executable); print(numpy.__version__)"
conda list numpy
conda deactivate
\end{tbash}


这时候应该显示 Python 3.10，解释器路径也应当落在刚创建的环境里，并且能够成功导入 NumPy。别只看提示符前面多了一个环境名，就认定切换完成了；我们再看一次实际路径、试一次导入，才能确认这条 python 命令真的用上了它。
\end{example}
\begin{note}
\texttt{conda remove --name 名称 --all} 会删除指定环境的全部包。先退出并核对名称；不要把平时使用的环境拿来练删除。环境解决软件依赖，ROS 2 的 Bash/zsh 加载脚本另按卷二说明，不把它们写入所有 Shell 的通用启动流程。
\end{note}
\begin{exercise}\label{ex:v1-conda-env}
独立完成 Python 3.10 与 NumPy 环境的创建、激活、导入验证和退出，记录解释器路径。确认环境只用于练习后删除，再检查环境列表与当前 Python 路径。保留日志，失败时记录下载或解算错误，不宣称“提示符变了就是成功”。
\end{exercise}
\begin{conclusion}
配置决定渠道与激活行为，环境决定这次项目用哪套解释器和包。先读实际值，再改用户范围；先验证路径与导入，再清理练习环境。ROS 2 工作空间使用其支持的系统环境，深度学习项目的 Conda 环境按需单独激活。
\end{conclusion}
